% !PRINTED 396
\weberpage{396}
Es ist also jeder der Complexe 6), 7) azygetisch mit dem Complexe 1),
und es muss also eine und nur eine der beiden Doppeltangenten
$\varrho, \varrho'$ in 6) vorkommen; ist dies $\varrho$, so ist $\varrho'$ in
7) enthalten. Ist $\varrho\sigma$ ein Paar von 6), so ist $\sigma$ von seinem
conjugirten $\sigma'$ verschieden, weil sonst $\sigma\varrho'$ in 7) und folglich
$\varrho\varrho'$ in 5) vorkommen müsste, was nicht der Fall ist. Also haben
wir, wenn $\tau, \tau'$ zwei weitere conjugirt imaginäre Doppeltangenten
sind, die Complexe:
\[
\begin{array}{rl}
6) & x\xi,\ p\eta,\ r\zeta,\ \varrho\sigma,\ \sigma'\tau \\
7) & x\xi',\ p\eta',\ r\zeta',\ \varrho'\sigma',\ \sigma\tau'.
\end{array}
\]
Ist $t\lambda$ das letzte Paar des Complexes 6), so kommt das Paar
$t'\lambda'$ in 7) vor, und diese beiden Paare müssen ein gemeinsames
Element enthalten. Da aber $\lambda$ nicht gleich $t'$ sein kann, weil sonst
beide Paare identisch wären, so muss $t = t'$ (oder $\lambda = \lambda'$, was
nicht wesentlich verschieden ist) sein; es ist also $t$ reell, und es
wird:
\[
\begin{array}{rl}
6) & x\xi,\ p\eta,\ r\zeta,\ \varrho\sigma,\ \sigma'\tau,\ t\lambda \\
7) & x\xi',\ p\eta',\ r\zeta',\ \varrho'\sigma',\ \sigma\tau',\ t\lambda'.
\end{array}
\]
Jetzt kehren wir zu dem syzygetischen Complextripel 1), 2), 3) zurück.
Da wir aus 6) schliessen, dass $x\,p\,t$ und $x\,r\,t$ azygetisch sind, so
kann $t$ weder in 2) noch in 3) vorkommen, und muss also in 1)
enthalten sein. Da aber 1) ein reeller Complex ist, so muss darin $t$
mit einer reellen Doppeltangente $u$ gepaart erscheinen, und der
Complex 1) wird
\[
x\,y,\ p\,q,\ r\,s,\ t\,u,
\]
wodurch unser Satz 3. bewiesen ist.

\gesperrt{4. Wenn in einem Complexe fünf reelle Paare vorkommen, so
sind alle Doppeltangenten reell.}

Gehen wir aus von einem Complexe mit fünf reellen Paaren
\[
1)\quad x_1y_1,\ x_2y_2,\ x_3y_3,\ x_4y_4,\ x_5y_5,
\]
und nehmen zunächst an, dass ausser dem letzten Paare dieses
Complexes noch eine imaginäre Doppeltangente $\xi$ existire, dann ist
auch eine conjugirte Doppeltangente $\xi'$ vorhanden, und der durch das
Paar $\xi\xi'$ bestimmte reelle Complex 2) ist mit 1) syzygetisch [weil
$\xi$ und $\xi'$ nicht in 1) vorkommen]. Die vier gemeinschaftlichen
Elemente von 1) und 2) müssen aber reell
%CONT>>
